% !PRINTED 662
\weberpage{662}
If we make this assumption about the functionals $\varphi_n$, then
we can apply \gesperrt{Kummer's} theorem in the form given in
\S.\,174, 3.

In what follows we denote by $\varepsilon$ unit functionals, of
which no closer knowledge is required, and set
\weq{7}{\Phi_n = \prod^{t} \varphi_{nt'}^{t} = \varepsilon\, \vartheta_n^m,}
where $t$ runs through the residues of $m$ that are not divisible by
$q$ and are less than $m$, and $t'$ is determined by
$t\,t' \equiv 1\ (\mathrm{mod}\ m)$. Then $\vartheta_n$ is a
cyclotomic number in a higher field, whose $m$th power is contained
in $\Omega_m$ and which satisfies the second condition, that
\weq{8}{\vartheta_n^{-1} \vartheta_1^n = \alpha}
is a number of the field $\Omega_m$.

We now introduce a symbol that is used in many places, by
understanding by $E(x)$ the greatest integer contained in the real
number $x$, and by $(x)$ the residue, which is always a proper
fraction and which, should $x$ itself be an integer, is to be set
equal to zero. Then
\weq{9}{x = E(x) + (x).}

If $x$ is an integer, then according to this mode of notation
\[
m \left(\frac{x}{m}\right)
\]
is likewise an integer, and indeed the least positive residue of the
division of $x$ by $m$.

In formula (7) the index of the functionals $\varphi$ is determined
only with respect to the modulus $m$. The exponent, however, is
reduced to its least residue with respect to the modulus $m$. If,
then, we replace the index $t'$ in (7) by $n' t'$, then $t$ must be
replaced by the residue of the division of $n t$ by $m$, that is
[by (9)], by
\[
m \left(\frac{n t}{m}\right) = n t - m E\left(\frac{n t}{m}\right)
\]
and we can set
\[
\Phi_n = \prod^{t} \varphi_{t'}^{\,n t - m E\left(\frac{n t}{m}\right)} = \varepsilon\, \vartheta_n^m,
\]
on the other hand
\[
\Phi_1 = \prod^{t} \varphi_{t'}^{t} = \varepsilon\, \vartheta_1^m,
\]
%CONT>>
